\chapter{Comments on Existence Theory for Parabolic PDE}

\section*{4.1 \quad Linear scalar PDE}

Consider a second order PDE on $\Omega \subset \mathbb R^n$ for a function
$u : \Omega \to \mathbb R$ of the form
\[
\frac{\partial u}{\partial t}
= a_{ij} \partial_i \partial_j u + b_i \partial_i u + cu
\tag{4.1.1}
\]
where $\partial_i := \frac{\partial}{\partial x^i}$ and $a_{ij}, b_i, c :
\Omega \to \mathbb R$ are smooth coefficients. Such an equation is called
\emph{parabolic} if $a_{ij}$ is uniformly positive definite, that is, if there
exists $\lambda > 0$ such that
\[
a_{ij}\xi_i\xi_j \ge \lambda\lvert \xi\rvert^2
\tag{4.1.2}
\]
for all $\xi \in \mathbb R^n$. The heat equation
\[
\frac{\partial u}{\partial t} = \Delta u
\]
is a parabolic PDE.

This notion extends to manifolds. Let $\M$ be closed and consider the PDE for
$u : \M \to \mathbb R$ given by
\[
\frac{\partial u}{\partial t} = L(u)
\tag{4.1.3}
\]
where $L : C^\infty(\M) \to C^\infty(\M)$ can be written with respect to local
coordinates $\{x^i\}$ as
\[
L(u) = a_{ij}\partial_i\partial_j u + b_i\partial_i u + cu
\]
for $a_{ij}$, $b_i$ and $c$ locally defined smooth, real coefficients. Equation
(4.1.3) is parabolic if $a_{ij}$ is positive definite for all $x \in \M$.
This is well defined, independent of the choice of coordinates. See also
Section~4.2.

There is a good theory for such equations. Given a smooth initial function
$u_0 : \M \to \mathbb R$, there exists a smooth solution $u : \M \times
[0,\infty) \to \mathbb R$ to
\[
\left\{
\begin{aligned}
\frac{\partial u}{\partial t} &= L(u) && \text{on } \M \times [0,\infty) \\
u(0) &= u_0 && \text{on } \M.
\end{aligned}
\right.
\]
There is also uniqueness: if $u$ and $v$ solve $\frac{\partial u}{\partial t}
= L(u)$ and $\frac{\partial v}{\partial t} = L(v)$ on $\M \times [0,T]$, and
either $u(0)=v(0)$ or $u(T)=v(T)$, then $u(t)=v(t)$ for all $t\in[0,T]$.

\section*{4.2 \quad The principal symbol}

\begin{definition}
\label{topping-pde-principal-symbol-scalar}
\dcref{4.2.1}
\group{topping-pde-existence}
\source{topping:page-0055}
The principal symbol $\sigma(L) : T^*\M \to \mathbb R$ of the operator $L$ is
defined by
\[
\sigma(L)(x,\xi) = a_{ij}(x)\xi_i\xi_j.
\]
\end{definition}

This is well defined, independent of the choice of coordinates. An alternative
definition is as follows: given $(x,\xi) \in T^*\M$, and smooth functions
$\phi,f : \M \to \mathbb R$ with $d\phi(x)=\xi$, define
\[
\sigma(L)(x,\xi)f(x) = \lim_{s \to \infty} s^{-2}e^{-s\phi(x)}
L(e^{s\phi}f)(x).
\tag{4.2.2}
\]

\begin{remark}
\label{topping-pde-principal-symbol-limit-justification}
\dcref{4.2.2}
\group{topping-pde-existence}
\source{topping:page-0056}
The two definitions of $\sigma(L)$ agree. Differentiating $e^{s\phi}f$ shows
that the leading term of $e^{-s\phi(x)}L(e^{s\phi}f)(x)$ is
$s^2 a_{ij}\partial_i\phi(x)\partial_j\phi(x)f(x)$, so the displayed limit is
$a_{ij}\xi_i\xi_j f(x)$.
\end{remark}

\begin{remark}
\label{topping-pde-principal-symbol-fourier-convention}
\dcref{4.2.1}
\group{topping-pde-existence}
\source{topping:page-0056}
Some authors write $i\xi$ where this text writes $\xi$, in accordance with the
Fourier transform convention.
\end{remark}

Thus the PDE (4.1.3) is parabolic if $\sigma(L)(x,\xi) > 0$ for all
$(x,\xi)\in T^*\M$ with $\xi\neq 0$.

\begin{example}
\label{topping-pde-laplace-beltrami-example}
\dcref{4.2.2}
\group{topping-pde-existence}
\source{topping:page-0056}
For the Laplace-Beltrami operator
\[
\Delta=\frac{1}{\sqrt{g}}\partial_i\bigl(\sqrt{g}\,g^{ij}\partial_j\bigr)
= g^{ij}\partial_i\partial_j + \text{lower order terms},
\]
where $g:=\det g_{ij}$, one has
\[
\sigma(\Delta)(x,\xi)=g^{ij}\xi_i\xi_j=|\xi|^2>0.
\]
Hence the heat equation $\frac{\partial u}{\partial t}=\Delta u$ is parabolic
in this setting also.
\end{example}

\section*{4.3 \quad Generalisation to vector bundles}

The theory of parabolic equations extends under various generalisations of the
concept of parabolic, including higher order equations, less regular
coefficients and initial data, time-dependent coefficients, vector bundle
sections, and nonlinear equations. The text now turns to the vector bundle
case.

Let $E$ be a smooth vector bundle over a closed manifold $\mathcal M$.
Whereas above, we considered functions $u : \mathcal M \to \mathbb R$, now we
consider sections $v \in \Gamma(E)$. Locally we may write $v = v^\alpha e_\alpha$
for some local frame $\{e_\alpha\}$. Consider the equation
\[
\frac{\partial v}{\partial t} = L(v)
\tag{4.3.1}
\]
where $L$ is a linear second order differential operator, meaning that $L :
\Gamma(E) \to \Gamma(E)$ may be given locally, in terms of local coordinates
$\{x^i\}$ on $\mathcal M$ and a local frame $\{e_\alpha\}$ on $E$, as
\[
L(v) = \bigl[a^{ij}_{\alpha\beta}\partial_i\partial_j v^\beta
  + b^i_{\alpha\beta}\partial_i v^\beta
  + c_{\alpha\beta}v^\beta\bigr] e_\alpha.
\]

The principal symbol now is a vector bundle homomorphism
$\sigma(L) : \Pi^*(E) \to \Pi^*(E)$, where $\Pi : T^*\mathcal M \to \mathcal M$
denotes the bundle projection. The fibre of $\Pi^*(E)$ over $(x,\xi)\in
T^*(\mathcal M)$ is $E_x$. The symbol is
\[
\sigma(L)(x,\xi)v = \bigl(a^{ij}_{\alpha\beta}\xi_i\xi_j v^\beta\bigr)e_\alpha.
\tag{4.3.2}
\]

Again, there is a coordinate-independent definition: given $(x,\xi)\in T^*M$,
for all $v\in \Gamma(E)$ and $\phi:M\to \mathbb R$ with $d\phi(x)=\xi$,
define
\[
\sigma(L)(x,\xi)v=\lim_{s\to\infty}s^{-2}e^{-s\phi(x)}L\bigl(e^{s\phi}v\bigr)(x).
\tag{4.3.3}
\]

For a vector bundle with connection $\nabla$, if $\Delta$ denotes the
connection Laplacian, then
\[
\Delta\bigl(e^{s\phi}v\bigr)=e^{s\phi}\bigl(s^2|d\phi|^2v+\text{lower order terms in }s\bigr),
\]
and so $\sigma(\Delta)(x,\xi)v=|\xi|^2v$ because $d\phi=\xi$. Thus
$\sigma(\Delta)(x,\xi)=|\xi|^2\mathrm{id}$.

We say that (4.3.1) is strictly parabolic if there exists $\lambda>0$ such that
\[
\langle \sigma(L)(x,\xi)v,v\rangle \ge \lambda |\xi|^2|v|^2
\tag{4.3.4}
\]
for all $(x,\xi)\in T^*M$ and $v\in \Gamma(E)$. This notion is also often
called strongly parabolic to distinguish it from more general definitions.

\begin{remark}
\label{topping-pde-vector-parabolic-fibre-metric}
\dcref{4.3.1}
\group{topping-pde-existence}
\source{topping:page-0058}
Equation (4.3.1) is parabolic if (4.3.4) holds for any fibre metric
$\langle\cdot,\cdot\rangle$ on $E$.
\end{remark}

\begin{remark}
\label{topping-pde-vector-parabolic-positive-multiple}
\dcref{4.3.2}
\group{topping-pde-existence}
\source{topping:page-0058}
On a closed manifold $M$, equation (4.3.1) is parabolic if, for all
$(x,\xi)\in T^*M$ with $\xi\ne 0$, the symbol $\sigma(L)(x,\xi)$ is a positive
multiple of the identity on $E_x$.
\end{remark}

In practice we need to discuss nonlinear PDE of the form
\[
\frac{\partial v}{\partial t}=P(v)
\tag{4.3.5}
\]
where $P$ is a quasilinear second order differential operator, meaning that
$P:\Gamma(E)\to\Gamma(E)$ may be given locally, in terms of local coordinates
$\{x^i\}$ on $M$ and a local frame $\{e_\alpha\}$ on $E$, as
\[
P(v)=\bigl[a^{ij}{}_{\alpha\beta}(x,v,\nabla v)\,\partial_i\partial_j v^\beta
+b^\alpha(x,v,\nabla v)\bigr]\,e_\alpha.
\]
The Ricci flow equation $\frac{\partial g}{\partial t}=-2\operatorname{Ric}(g)$
fits into this category.

Given an arbitrary fixed section $w\in\Gamma(E)$ independent of $t$, we call
(4.3.5) parabolic at $w$ if the linearisation of (4.3.5) at $w$,
\[
\frac{\partial v}{\partial t}=[DP(w)]v,
\]
is parabolic as above.

\section*{4.4 \quad Properties of parabolic equations}

If $\frac{\partial v}{\partial t}=P(v)$ is parabolic at $w$ and the
$a_{\alpha\beta}^{ij}$ and lower order terms satisfy suitable regularity
hypotheses, then there exist $\varepsilon>0$ and a smooth family $v(t)\in
\Gamma(E)$ for $t\in[0,\varepsilon]$ such that
\[
\left\{
\begin{aligned}
\frac{\partial v}{\partial t} &= P(v) \qquad t\in[0,\varepsilon] \\
v(0) &= w.
\end{aligned}
\right.
\]

There is also uniqueness: if $v$ and $w$ solve $\frac{\partial v}{\partial t}
= P(v)$ and $\frac{\partial w}{\partial t}=P(w)$ for $t\in[0,\varepsilon]$,
and either $v(0)=w(0)$ or $v(\varepsilon)=w(\varepsilon)$, then
$v(t)=w(t)$ for all $t\in[0,\varepsilon]$.
